\documentclass[11pt,a4paper]{article}

\usepackage[UTF8]{ctex}
\usepackage{amsmath,amssymb,bm}
\usepackage{enumitem}
\usepackage{geometry}
\usepackage{fancyhdr}
\usepackage{xcolor}

\geometry{left=2.2cm,right=2.2cm,top=2.1cm,bottom=2.1cm}
\setlength{\headheight}{14pt}
\setlength{\parindent}{2em}
\setlength{\parskip}{0.25em}
\setlist[enumerate,1]{label=\textbf{(\alph*)},leftmargin=2.7em,itemsep=0.45em,topsep=0.35em}

\newcommand{\R}{\mathbb{R}}
\newcommand{\E}{\mathbb{E}}
\newcommand{\Var}{\operatorname{Var}}
\newcommand{\Cov}{\operatorname{Cov}}
\newcommand{\grad}{\nabla}
\newcommand{\points}[1]{\hfill\textcolor{gray}{[#1 分]}}

\pagestyle{fancy}
\fancyhf{}
\lhead{AI3002 人工智能与机器学习基础}
\rhead{HW0：数学基础}
\cfoot{\thepage}
\renewcommand{\headrulewidth}{0.4pt}
\fancypagestyle{plain}{%
  \fancyhf{}%
  \lhead{AI3002 人工智能与机器学习基础}%
  \rhead{HW0：数学基础}%
  \cfoot{\thepage}%
  \renewcommand{\headrulewidth}{0.4pt}%
}

\title{\vspace{-1.4cm}\textbf{人工智能与机器学习基础\\HW0：数学基础诊断}}
\author{中国科学技术大学}
\date{2026 年秋季学期\quad 发布日期与截止日期待定}

\begin{document}

\maketitle
\pagestyle{fancy}
\thispagestyle{fancy}

\section*{说明}

本次作业用于帮助大家检查学习本课程所需的数学基础，内容限于线性代数、微积分、概率统计和基础优化。预计完成时间为 \textbf{2--3 小时}。

请写出必要的计算过程；仅给出最终数值可能无法获得全部分数。可以查阅教材和课程提供的复习资料，但请独立完成。除简单算术核对外，不需要使用编程或符号计算软件。

本作业满分 100 分，最终只记录完成状态，不计入正式书面作业平均分。

提交一个 PDF 文件，命名为：
\[
    \texttt{学号-姓名-HW0.pdf}.
\]

\section{向量与矩阵\points{20}}

设数据矩阵 $X\in\R^{n\times d}$，参数向量 $\bm w\in\R^d$，标签向量 $\bm y\in\R^n$，并定义预测向量
\[
    \hat{\bm y}=X\bm w.
\]

\begin{enumerate}
    \item 分别写出 $X^\top X$、$X^\top\bm y$、$XX^\top$ 和 $\hat{\bm y}-\bm y$ 的维度。\points{6}

    \item 取
    \[
        X=\begin{bmatrix}
            1&2\\
            0&-1\\
            2&1
        \end{bmatrix},\qquad
        \bm w=\begin{bmatrix}1\\-1\end{bmatrix},\qquad
        \bm y=\begin{bmatrix}0\\1\\2\end{bmatrix}.
    \]
    计算 $X\bm w$、$X\bm w-\bm y$ 以及 $X^\top(X\bm w-\bm y)$。\points{8}

    \item 对向量 $\bm a=(1,2,-1)^\top$ 和 $\bm b=(2,0,2)^\top$，计算它们的内积、$\ell_2$ 范数以及余弦相似度
    \[
        \frac{\bm a^\top\bm b}{\lVert\bm a\rVert_2\lVert\bm b\rVert_2}.
    \]
    根据结果说明两向量是否正交。\points{6}
\end{enumerate}

\section{特征值与二次型\points{15}}

给定对称矩阵
\[
    A=\begin{bmatrix}2&1\\1&2\end{bmatrix}.
\]

\begin{enumerate}
    \item 求 $A$ 的两个特征值，并为每个特征值写出一个对应的单位特征向量。\points{8}

    \item 判断 $A$ 是否为正定矩阵，并说明理由。你可以使用特征值，也可以直接整理二次型
    \[
        \bm x^\top A\bm x,\qquad \bm x=(x_1,x_2)^\top.
    \]
    \points{5}

    \item 不显式计算矩阵逆，利用特征向量或直接求解线性方程，求满足
    \[
        A\bm x=\begin{bmatrix}3\\3\end{bmatrix}
    \]
    的 $\bm x$。\points{2}
\end{enumerate}

\section{多元微积分与链式法则\points{25}}

设 $W\in\R^{m\times d}$、$\bm x\in\R^d$、$\bm b,\bm y\in\R^m$，定义
\[
    \bm z=W\bm x+\bm b,
    \qquad
    L=\frac{1}{2}\lVert\bm z-\bm y\rVert_2^2.
\]

\begin{enumerate}
    \item 将 $L$ 写成求和形式，即写成关于 $z_i$ 和 $y_i$ 的表达式，并求 $\partial L/\partial z_i$。\points{5}

    \item 使用链式法则求 $\grad_{\bm b}L$ 和 $\grad_{\bm x}L$。请同时写出它们的维度。\points{8}

    \item 求 $\grad_W L$。答案应为一个 $m\times d$ 矩阵，并尽量写成向量外积的形式。\points{8}

    \item Sigmoid 函数定义为
    \[
        \sigma(t)=\frac{1}{1+e^{-t}}.
    \]
    证明
    \[
        \sigma'(t)=\sigma(t)\bigl(1-\sigma(t)\bigr).
    \]
    \points{4}
\end{enumerate}

\section{概率与统计\points{25}}

\begin{enumerate}
    \item 某事件 $D$ 的先验概率为 $P(D)=0.1$。一个检测方法满足
    \[
        P(+\mid D)=0.9,
        \qquad
        P(-\mid D^c)=0.95.
    \]
    计算 $P(+)$ 和 $P(D\mid +)$。请写出使用的全概率公式和贝叶斯公式。\points{10}

    \item 离散随机变量 $X$ 的分布如下：
    \[
    \begin{array}{c|ccc}
        x & -1 & 0 & 2\\ \hline
        P(X=x) & 0.2 & 0.5 & 0.3
    \end{array}
    \]
    计算 $\E[X]$、$\E[X^2]$ 和 $\Var(X)$。\points{9}

    \item 判断下面的说法是否正确，并用一句话说明理由：
    \begin{quote}
        若 $\Cov(X,Y)=0$，则随机变量 $X$ 与 $Y$ 一定相互独立。
    \end{quote}
    如果你认为不正确，请给出一个简单反例。\points{6}
\end{enumerate}

\section{梯度下降\points{15}}

考虑函数
\[
    f(w_1,w_2)=(w_1-1)^2+2(w_2+1)^2.
\]

\begin{enumerate}
    \item 求梯度 $\grad f(w_1,w_2)$，并写出函数的全局最小点。\points{5}

    \item 从 $\bm w^{(0)}=(0,0)^\top$ 出发，使用学习率 $\eta=0.25$ 做一步梯度下降：
    \[
        \bm w^{(1)}=\bm w^{(0)}-\eta\grad f\bigl(\bm w^{(0)}\bigr).
    \]
    求 $\bm w^{(1)}$ 以及 $f(\bm w^{(0)})$、$f(\bm w^{(1)})$。\points{6}

    \item 若改用学习率 $\eta=1$，同样从 $\bm w^{(0)}$ 做一步更新，目标函数会下降还是上升？这个例子说明学习率过大可能产生什么问题？\points{4}
\end{enumerate}

\section*{完成情况反馈（不计分）}

请在提交末尾简要填写：
\begin{enumerate}[label=\arabic*.,leftmargin=2em]
    \item 完成本次作业大约花费了多长时间？
    \item 哪一部分最不熟悉：线性代数、微积分、概率统计，还是优化？
    \item 是否希望课程提供某一部分的补充讲义或习题课？
\end{enumerate}

\vfill
\begin{center}
    \textcolor{gray}{--- End of HW0 ---}
\end{center}

\end{document}
